\section{L'angle d'holonomie $\theta_p$ comme objet primitif}
\label{sec:holonomie_primitive}

Cette section établit le résultat fondamental qui distingue la
PT des théories classiques de jauge : sous les hypothèses
structurales de la PT, l'angle d'holonomie modulaire $\theta_p$
est l'\emph{unique} invariant primitif satisfaisant
simultanément quatre conditions naturelles, et le couple
classique $(A_\mu, F_{\mu\nu})$ est entièrement reconstructible
comme dérivée d'un objet plus primitif --- le potentiel scalaire
de persistance $S = -\ln \alpha$. La position C du débat
fondationnel n'est plus une interprétation : elle est une
conséquence structurale.

\subsection{Proposition 3.1 : $A_\mu$ et $F_{\mu\nu}$ sont dérivés}
\label{ssec:prop_3_1}

\begin{proposition}[Décomposition du potentiel de persistance]
\label{prop:A_F_derives}
Soit $S(\mu) = -\ln \alpha(\mu)$ le potentiel de persistance. Alors
\begin{itemize}
  \item le 4-potentiel $A_\mu$ est identifiable au gradient
  $\partial_\mu S$ ;
  \item le tenseur de champ $F_{\mu\nu}$ est identifiable aux
  composantes antisymétriques de $\mathrm{Hess}(S)$ ;
  \item la métrique $g_{\mu\nu}$ est identifiable aux composantes
  symétriques de $\mathrm{Hess}(S)$ (Lemme F, ch12 monographie).
\end{itemize}
\end{proposition}

\begin{proof}[Démonstration]
La construction est détaillée au chapitre~12 de la
monographie~\cite{PT_Monograph} ; on en retrace ici les trois
étapes clés.

\textbf{Étape 1 (identification du gradient).}
La métrique de Fisher associée à la distribution stationnaire
$\pi_m$ du crible se calcule comme le Hessien du logarithme de la
fonction de partition $Z_{\rm Ruelle}$, et la fonction
$\alpha(\mu)$ apparaît comme la première fonction propre du
transfert. Par construction, le gradient $\partial_\mu S =
-\partial_\mu \ln \alpha$ joue dans le formalisme PT le rôle que
joue $A_\mu$ dans l'électrodynamique classique : il intervient
linéairement dans le terme d'interaction du Lagrangien effectif
(\S~F3 de la monographie) et il subit, sous les transformations
de jauge modulaire $\theta_p \mapsto \theta_p + 2\pi k$, un
déphasage compatible avec l'invariance de jauge classique.

\textbf{Étape 2 (identification du Hessien antisymétrique).}
Les composantes antisymétriques du Hessien
$\mathrm{Hess}(S)_{\mu\nu} = \partial_\mu \partial_\nu S$ encodent
les courbures dérivées du potentiel scalaire, et reproduisent les
six degrés de liberté de $F_{\mu\nu} = \partial_\mu A_\nu -
\partial_\nu A_\mu$ par identification directe. L'identité
$dA = F$ devient, sous PT, l'identité tautologique
$d(dS) = 0$ projetée sur les composantes antisymétriques de
$\mathrm{Hess}(S)$. Cette projection est non-triviale en présence
de la signature lorentzienne $g_{00} < 0$, qui sépare les
composantes spatiales et temporelles du Hessien.

\textbf{Étape 3 (identification du Hessien symétrique).}
Les composantes symétriques de $\mathrm{Hess}(S)$ reproduisent la
métrique $g_{\mu\nu}$ de l'espace-temps émergent (Bianchi I sur
les premiers actifs $\{3, 5, 7\}$, cf.\ Lemme F du
ch.~12~\cite{PT_Monograph}). Cette identification est, en
elle-même, le contenu du théorème de reconstruction métrique :
la métrique de l'espace-temps de la PT \emph{est} le Hessien du
logarithme de la fonction de partition, sans nouvelle structure
introduite.
\end{proof}

La conséquence immédiate est que $A_\mu, F_{\mu\nu}, g_{\mu\nu}$
sont tous trois \emph{dérivables du même potentiel scalaire}
$S(\mu)$, et que ce potentiel est, en retour, calculé à partir du
seul champ $\{g_n\}$. Aucun de ces trois objets n'est donc
primaire dans l'architecture PT : ils sont tous trois secondaires
par rapport à $S$, et $S$ est lui-même secondaire par rapport à
$\{g_n\}$. C'est sur ce constat structural que repose
l'argument d'unicité de la sous-section suivante.

\subsection{Proposition 3.2 : unicité de $\theta_p$ comme invariant primitif}
\label{ssec:prop_3_2}

\begin{proposition}[Primauté holonomique]
\label{prop:theta_p_unique}
L'angle d'holonomie $\theta_p$ est l'unique invariant primitif
simultanément :
\begin{enumerate}[label=(\roman*),nosep]
  \item jauge-invariant,
  \item sensible à Aharonov--Bohm,
  \item borné,
  \item PT-définissable depuis BA0 seul.
\end{enumerate}
\end{proposition}

\begin{proof}[Démonstration]
Nous montrons par exclusion qu'aucun autre invariant naturel ne
satisfait simultanément les quatre conditions.

\textbf{Exclusion de $A_\mu$.}
Le potentiel $A_\mu$ viole la condition (i) : il dépend de la
jauge par construction. Aucune projection naturelle de $A_\mu$
n'est à la fois jauge-invariante et bornée tout en restant
sensible à Aharonov--Bohm dans la région où $F = 0$. La position
A du débat fondationnel est donc structurellement incompatible
avec les quatre conditions.

\textbf{Exclusion de $F_{\mu\nu}$.}
Le tenseur de champ $F_{\mu\nu}$ satisfait (i), (iii), (iv) mais
viole la condition (ii) : il est, par construction, identiquement
nul dans la région intérieure d'un trajet Aharonov--Bohm idéal,
et ne peut donc pas rendre compte du déphasage observé. La
position B du débat doit ajouter une structure non-locale
supplémentaire pour récupérer la sensibilité à AB ; cette
structure ajoutée n'est pas PT-définissable depuis BA0 seul.

\textbf{Exclusion des fonctions arbitraires de $\{g_n\}$.}
Une fonction arbitraire $F[\{g_n\}]$ peut satisfaire (i), (ii),
(iv) mais viole en général (iii) : sans structure de transport
parallèle modulaire, la fonction n'est pas naturellement bornée
et n'admet pas d'identification simple avec un invariant
topologique mesurable. Seules les fonctions qui se réduisent à
l'angle modulaire $\theta_p$ d'un canal CRT
$\mathbb{Z}/p\mathbb{Z}$ remplissent les quatre conditions
simultanément.

\textbf{Existence et unicité.}
L'angle $\theta_p$ défini par
l'identité~\eqref{eq:sin2_theta_fr} satisfait (i) (il est
jauge-invariant car défini comme intégrale fermée), (ii) (il est
non-nul dans toute région où la connexion modulaire est
non-triviale, y compris dans la région intérieure d'un cycle
AB), (iii) (il prend ses valeurs dans $[0, \pi]$ par
construction), et (iv) (il est défini par
l'identité~\eqref{eq:sin2_theta_fr}, qui ne fait intervenir que
$\{g_n\}$ via $q(\mu)$). C'est l'unique invariant primitif qui
réalise les quatre conditions.
\end{proof}

\subsection{La position Wu--Yang--Healey promue en théorème PT}
\label{ssec:wu_yang_promu}

Les Propositions~\ref{prop:A_F_derives}
et~\ref{prop:theta_p_unique} ont une lecture commune : la
position C du débat fondationnel
(\S~\ref{ssec:position_C}) --- la primauté de l'holonomie ---
n'est plus, dans le cadre PT, un choix d'interprétation parmi
d'autres. Elle est une conséquence structurale : sous les axiomes
de la PT, $\theta_p$ est l'\emph{unique} invariant primitif
disponible, et $A_\mu, F_{\mu\nu}, g_{\mu\nu}$ sont tous trois
dérivés de ce dernier via le potentiel scalaire $S$.

Ce statut a deux conséquences immédiates. D'une part, la PT
\emph{prédit} que toute mesure de $A_\mu$ ou $F_{\mu\nu}$ se
réduit, dans son contenu informationnel, à une mesure de $\theta_p$
sur un canal modulaire actif. Aucune information physique n'est
perdue par cette réduction --- la phase AB est intégralement
encodée dans $\theta_3$ (cf.\ \S~\ref{ssec:ab_consequence}). D'autre
part, la PT \emph{exclut} les ontologies de la position A (la
réalité du potentiel) et de la position B (la non-localité
intrinsèque du champ) : ces ontologies postulent des objets
primitifs qui, dans le cadre PT, ne peuvent exister
indépendamment du potentiel scalaire $S$ d'où elles sont dérivées.

La position C devient ainsi, sous PT, plus qu'une interprétation :
elle est le seul cadre dans lequel les quatre conditions
naturelles (i)--(iv) de la
Proposition~\ref{prop:theta_p_unique} peuvent être satisfaites
simultanément. Le débat fondationnel est, dans cette mesure
précise, tranché.
